\section{引言}

半导体器件设计与制造的核心问题之一是如何精确控制半导体内的杂质分布，以满足不同器件对电学特性的要求。杂质浓度及其分布直接影响pn结的击穿电压、势垒电容、载流子迁移率等关键参数，因此准确测量杂质分布是半导体材料与器件研究中的基本任务之一。

传统的杂质分布测量方法包括四探针法和霍耳效应逐次去层测量技术。这些方法通过逐层剥离样品并测量薄层电导或薄层霍耳电压，可以获得杂质浓度及其分布。然而，逐次去层测量操作繁琐、耗时长，且对样品具有破坏性。相比之下，电容-电压法（C--V法）利用pn结势垒电容随反向偏压变化的特性，通过测量不同直流偏压下的势垒电容，可以在不破坏器件的情况下快速获得单边突变pn结轻掺杂一侧的杂质浓度及其分布。该方法尤其适用于检测扩散工艺或离子注入工艺所引入的杂质分布。

本实验基于C--V法测量$\mathrm{p}^+\mathrm{n}$单边突变结轻掺杂一侧的杂质浓度分布。实验采用锁定放大器作为微弱信号检测的核心仪器，将pn结势垒电容转换为交流电压信号进行精确测量。通过对多个二极管样品的C--V特性进行系统测量，结合理论模型计算得到各二极管的杂质浓度分布$N(w)$，并通过$1/C^{2}$--$V_{\mathrm{R}}$曲线的直线段外推获得接触电势差$V_{\mathrm{D}}$。此外，实验还考察了锁定放大器在不同时间常量下抑制噪声的能力，以及二极管相位特性与反向偏压的关系。
